\subsection{The global fixed-label pair count}
\label{sec:flag-global-pair-count}

The original fixed-row inequality and embedded-isometry
genus mean have compatible masses. We first identify their
normalizations, and then perform the finite parameter sum.
All lattices and the modulus in this section are those of
the given full row system.

\begin{lemma}[The primitive last-vector row mass]
\label{lem:flag-row-genus-mass}
Let \(Q=X^\perp/X\) be the original isotropic quotient,
\(z\) the image of the last vector of norm \(-T\),
and \(W=z^\perp\subset Q\), with its induced positive
definite hermitian form. For a row \(\Lambda\) let
\(L_\Lambda\) be the integral lattice in \(W\) obtained
from the row quotient and intersection with \(W\).
For any isometry-invariant nonnegative function \(\phi\)
of these definite lattices, and any genus \(\mathcal G\)
represented by a row, one has
\[
 \sum_{\Lambda:\,L_\Lambda\in\mathcal G}
     w_{\rm row}(\Lambda)\phi(L_\Lambda)
  =(T-1)^f
    \sum_{[L]\in\mathcal G/\cong}
       \frac{\phi(L)}{|\operatorname{Aut}(L)|}.
 \tag{RC1}
\]
Here \(T\) is the original completely split rational
prime and \(f=[E:\mathbf Q]/2\).
\end{lemma}
\begin{proof}
The fixed isotropic coefficient ideals are \(A_a\).
By the given split row decomposition, the row quotient
is unimodular, \(z\) is primitive (its saturated line
lattice is exactly \(\mathcal O_Ez\)), and its orthogonal
complement has discriminant module isomorphic to
\(\mathcal O_E/T\mathcal O_E\). The last assertion
also follows from the primitive gluing argument below.

We check that the rows in the given full system are
precisely representatives of the unimodular lattices in
\(Q\) containing \(z\) primitively, up to isometries
fixing \(z\). In the original rational hyperbolic
decomposition, any such lattice \(M\) lifts to
\[
 \bigoplus_a A_ae_a\;\oplus M\;\oplus
 \bigoplus_a A_a^{-1}f_a
\]
where \(f_a\) are the original rational isotropic duals.
This is a full self-dual lattice because the displayed
hyperbolic module and \(M\) are self-dual. It belongs
to the original genus. Indeed at each finite place
of the totally real field the CM algebra is unramified
or split. In the split case a self-dual lattice is the
sum of a free lattice and its perfect dual; in the
unramified field case a self-dual hermitian lattice
has an orthogonal integral basis of unit norms. To
construct it, the nondegenerate residue hermitian form
has a nonzero norm vector: otherwise polarization by
scalars and the nonzero residue trace would make
every pairing zero. Split off a unit-norm line and
induct. Norms of units surject onto the units of the
base local field, by the residue-field norm and
successive lifting using the surjective residue trace.
Thus all self-dual lattices in this local space have
the same local isometry type, including at residue
characteristic two. This proves the required genus
membership. The fixed flag and its coefficient ideals
are unchanged, and primitivity of \(z\) is exactly the
last-vector condition. In particular the original
fixed-congruence conditions hold at the same primes:
the isotropic part of the lifted lattice is the
original fixed coefficient module. The completeness
clause of the full row system therefore supplies a row
representative.

Conversely an isometry of two row quotients fixing \(z\)
lifts to an isometry of the actual split row lattices
preserving the isotropic flag and fixing the last vector.
Use the rational isotropic dual bases in their respective
split decompositions, map the two hyperbolic bases onto
one another, and use the given quotient isometry on the
orthogonal summands. Both orthogonal summands identify
isometrically with \(Q\), also on their integral lattices,
since the split hyperbolic lattice is self-dual. The
last vector belongs to the orthogonal summand in each
decomposition. The lifted isometry therefore satisfies
the original row equivalence, and uniqueness of the
given representatives makes the rows equal. The
source row weight is, by its definition, the inverse
order of the isometry group of the quotient lattice
fixing \(z\). This is not the inverse order of the
whole indefinite lattice stabilizer.

Now put \(Z=\mathcal O_E z\). Unimodular lattices
\(M\subset Q=W\perp Ez\) with \(M\cap W=L\),
\(M\cap Ez=Z\) are classified by graphs of
anti-isometries between the discriminant modules of
\(L,Z\). To see this, \(M/(L\oplus Z)\) is a self-annihilating
submodule of \(L^\#/L\oplus Z^\#/Z\) for the induced
hermitian pairing modulo \(\mathcal O_E\). Primitivity
makes the two projections injective. Self-annihilation
then makes both surjective: for example a vector in the
annihilator of the first projection would give an element
of \(M/(L\oplus Z)\) supported in the first factor and
hence zero. The two discriminant modules are thus
isomorphic and the graph condition is exactly the
negative-pairing condition. Conversely every such graph
gives a self-dual lattice with both intersections unchanged.

Since \(Z^\#/Z\simeq \mathcal O_E/T\mathcal O_E\),
the definite lattices in \(\mathcal G\) have this
same cyclic discriminant module. The algebra at \(T\)
is a product of \(f\) pairs of copies of \(\mathbf F_T\),
with conjugation exchanging the factors of each pair.
For a generator of a cyclic perfect hermitian module,
its pairing (after multiplication by \(T\)) is a
unit fixed by conjugation. Thus the possible
anti-isometries are given on each pair by two nonzero
scalars with prescribed nonzero product. There are
exactly \((T-1)^f\) such anti-isometries. This holds also
when \(T=2\). The group \(\operatorname{Aut}(L)\)
acts on these graphs, and the stabilizer of a graph
is the automorphism group of the glued lattice fixing
\(z\) (an isometry fixing \(z\) acts identically on \(Ez\)).
For a finite group acting on a finite set, the sum over
orbits of inverse stabilizer orders equals the cardinality
of the set divided by the group order. Applying this
identity to the graphs for each isometry class of \(L\)
proves (RC1), with exactly the original row weights.
\end{proof}

In particular, if an isometry-invariant nonnegative
quantity has genus sum at most \(\operatorname{Mass}(L)B\),
(RC1) with this quantity and with \(1\) bounds its
row-weighted sum by \(\operatorname{massRows} B\).
One applies the assertion within each genus and then
adds them. This is the source's \(\mathrm{RowsMeanC}\);
there is no factor depending on the number of genera
or on \(T\) in the resulting inequality.

\begin{theorem}[The original global fixed-label bound]
\label{thm:flag-global-pair-count}
In the original \(\mathrm{KTwoHypotheses}\), fix
\(C_H\ge0\), put \(r=r_S\), and take
\[
 C_a=\max(1,2(C_H+1),C_T),\quad
 C_h=300000 C_a,\quad C_p=16.
\]
The choices \(C'_0=1,\ C_{l,0}=42C_a\) suffice for
\(\mathrm{PairsP}\). In particular, for any finite set
\(\mathscr S\) of actual compatible relevant zero-Gram
term pairs with fixed positive indices and labels
\((j,\delta),(j',\delta')\), writing
\(\Delta=k2Del(j,\delta)\), \(S=k2S(H,j,\delta)\)
and similarly \(\Delta',S'\), with \(H=b^{C_Hm}\),
one has
\[
 \sum_{(t,t')\in\mathscr S}
   \mathrm{wt}(t)\mathrm{wt}(t')\,\omega(t)\omega(t')
 \le \operatorname{massRows}\,
       E(\Delta,S) E(\Delta',S'),
 \tag{RC2}
\]
where
\[
 E(\Delta,S)=\Delta^{k/2}
     \bigl(b^{C_hm^2 f}\bigr)^{C_hr^4}S^{16k}.
\]
Thus the factor is exactly the original \(k2e\), with
its original exponent \(C_p k\).
\end{theorem}
\begin{proof}
For the chosen original data, let \(\rho=\operatorname{rd}(E)\),
\(g=g_S\), and \(T=T_u\). The size conditions give
\[
 \begin{gathered}
 m=2r+2,\quad 42C_a\le m,\quad
 m^4\le f,\quad m^5 f\le k,\quad
 42C_a m^2\log b\le k,\quad \log f\le\sqrt r,\\
 2\le\rho\le b,\qquad (kb)^8\le g,\qquad
 1\le T\le b^{C_am}.
 \tag{RC3}
 \end{gathered}
\]
Indeed the original scale satisfies \(\log b=L>1\).
The lower bound \(C'm^5 fL\le k\) and \(C'\ge1\)
give the displayed rank bound; the \(C_l\) bounds give
the two displayed bounds with \(42C_a\). The bound
\(\log\rho\le L\) follows from the root-discriminant
condition with its positive denominator at least one
in the admitted size range. The original \(g_S\) is
the product of the two positive level factors, and the
first is at least \((kb)^8\). The original prime
selection gives the completely split \(T\) and its
stated upper bound. These are the original
\(\mathrm{pairsPar}\) inequalities.

If \(\mathscr S\) is empty there is nothing to prove.
Otherwise all parameters for both labels occur on
actual terms. The original normalized leading minor
of order \(l\) has norm
\[
 N_l=|\mathrm N_{E/\mathbf Q}\delta_l|\,
        \mathrm N(\mathfrak d)^{-2\operatorname{frameExp}(r,l)}
    =\prod_{\tau} a_{\tau,l}^{\,2},\qquad
 a_{\tau,l}=|\tau(\delta_l)|\rho^{-2\operatorname{frameExp}(r,l)}.
 \tag{RC4}
\]
The equality follows from the CM product formula
and \(\mathrm N(\mathfrak d)=\rho^{2f}\), with one
embedding from each complex pair. The transported
integral leading ideals \(A_a=\mathfrak d^{-(r-a)}\)
show that the fractional ideal of which \(N_l\) is the
norm is integral: each cross pairing on \(A_ae_a,A_bf_b\)
is integral, so multiplying the determinant by
\((\prod_{a<l}A_a)^2\) gives an integral ideal.
Thus every nonzero \(N_l\) is at least \(1\).
The normal-form label theorem gives
\(\Delta=N_j>0\) and \(D_b=\prod_{l\in\mathcal E}N_l\)
for the original nonzero-minor set \(\mathcal E\).
The source height factor is
\[
 S=\prod_\tau\max\bigl(1,\max_{l\in\mathcal E}
              a_{\tau,l}/(2H)^{2r}\bigr).
\]
Consequently \(S\ge1\) and
\(N_l\le (2H)^{4rf} S^2\) for every \(l\in\mathcal E\).
Since \(4rf(1+C_Hm)\le2(C_H+1)m^2f\), put
\[
 A=b^{C_a m^2 f},\qquad U=AS^2,\qquad U'=AS'^2.
\]
We have \(1\le\Delta\le U,\ 1\le D_b\le U^r\)
and the analogous bounds for the second label;
also \(N_T=T^f\le A\).

Let \(P(x,x')\) be the explicit \(\mathrm{pairBd}\)
of Theorem~\ref{thm:embedded-isometry-genus-mean}
with the original rank parameter \(r\), and
\(d(x),d'(x')\) the \(\mathrm{datBd}\) factors of
Theorem~\ref{thm:fixed-normal-form-weighted-sum}.
By the original label congruence both terms in a
compatible pair have equal target and source rows.
Therefore their two square-root row weights multiply
to the product of the target and source row weights.
Partition \(\mathscr S\) by its target row, use
Theorem~\ref{thm:weighted-fixed-row-comparison}, and
multiply each resulting estimate by that target row
weight. Summing over the actual row representatives
and interchanging the finite index sums gives
\[
 \sum_{\mathscr S}\mathrm{wt}(t)\mathrm{wt}(t')
             \omega(t)\omega(t')
 \le \operatorname{massRows}
        \sum_{1\le x\le\lfloor\Delta\rfloor}
        \sum_{1\le x'\le\lfloor\Delta'\rfloor}
              P(x,x')d(x)d'(x').
 \tag{RC5}
\]
Indeed the embedded triples are counted by an
isometry-invariant function of the definite lattice
(conjugate the two embedded data and their comparison
isometry). Its genus sum has the required mass factor
by the proved embedded-isometry mean. The last-vector
quotient has rank \(k\) and the original discriminant
module, by the preceding lemma. Applying (RC1) to
each such function yields (RC5). The genus bound is
finite; all summands are nonnegative. The argument
also applies to the source's extended nonnegative-real
count, followed by restriction to finite real bounds.

It remains to bound this actual finite sum.
Set \(c=1+\log_2(2f)\le4r\),
\(D_0=T^{2f}\), \(Y=xx'D_0^2\), and
\(M=(20f\rho^2)^{2f(6+c)}\rho^f\).
Since \(D_0\le A^2,\ x\le U,\ x'\le U'\),
\[
 Y\le A^6(SS')^2\le (UU')^3 .
\]
The size bounds imply \(f\le b^{2r}\),
\(8^{r+1}\le A,\ f^2\le A\).
Moreover \(M\le b^{181fr^2}\): use
\(20f\rho^2\le b^{2r+7}\),
\(2(2r+7)(6+4r)+1\le181r^2\).
Thus \(2+2M\le A^{183}\).
Define
\[
 q_0=(32/g)^{kf}(4kf)^{4f}A^{42}.
\]
From (RC3), \(A^{42}\le e^{kf}\);
\(f\le k,\ b\ge2,\ k\ge9\) give
\[
 q_0^{1/f}\le (32e/g)^k (4kf)^4
       \le 256\,2^{-k} k^{8-8k}\le\tfrac12 .
\]
Hence \(q_0\le1/2\).
Writing \(Z_0=(SS')^{14}\ge1\), the two
bases of powers in the bracket of \(\mathrm{pairBd}\)
are bounded by \(Z_0/2\): the first is
\((32/g)^{kf}(4kf)^{4f}Y^7\), and the second is no
larger because \(D_0xx'\le Y^7\).
The geometric sum \(\sum_{s=1}^{k-1}2^{-s}\le1\)
therefore bounds the entire bracket by
\((2+2M)Z_0^k\).
As \(6r+10+4c\le32r\), it follows that
\[
 P(x,x')\le (xx')^{k/2}
       (UU')^{96r+183}(SS')^{14k}.
 \tag{RC6}
\]

The exact factor \(\mathrm{datBd}\) is
\[
 d(x)=(\Delta/x)^{k/2}
 8^{r+1}f^2\Delta^{3c+6}N_T^2
             (D_b\Delta N_T)^{10000(r+1)^3}.
\]
Its factor after \((\Delta/x)^{k/2}\) is at most
\(U^{250000r^4}\): the first two factors cost
\(U^2\), the next powers cost \(U^{12r+8}\),
and the last costs at most \(U^{240000r^4}\).
The analogous bound holds for \(d'(x')\).
In (RC6) the powers of \(x,x'\) thus cancel exactly.
There are at most \(UU'\) choices of the two indices.
Consequently the sum in (RC5) is bounded by
\[
 \Delta^{k/2}\Delta'^{k/2}
       (UU')^{300000r^4}(SS')^{14k},
\]
since \(250000r^4+96r+184\le300000r^4\).
Finally \(k\ge m^5 f\ge m^9\) and \(m\ge42\)
imply \(300000r^4\le k\). Hence
\[
 U^{300000r^4}S^{14k}
 \le \bigl(b^{C_h m^2 f}\bigr)^{C_h r^4}S^{16k},
\]
and likewise for the other label. This proves (RC2)
with constants chosen before the arithmetic data vary.
\end{proof}

This is the fixed-label global arithmetic supplier
\(\mathrm{PairsP}\), not the full zero-Gram return
estimate. The latter must still combine (RC2) with
the actual amplitude and distinct-label estimates,
and then handle generic returns and the finite columns.
